%! AP Statistics — Unit 4, Lesson 4.3 — Group Activity KEY
\documentclass[10pt]{article}
\usepackage{apstats-article}
\usepackage{apstats-key}

\begin{document}

\pageheader{Unit 4, Lesson 4.3}{Group Activity --- KEY}
\namepartnerperiod

\vspace{0.1in}

\begin{scenariobox}[Context]{sky}
\small
A standard deck of 52 playing cards has 4 suits (hearts, diamonds, clubs, spades),
13 cards per suit (Ace, 2--10, Jack, Queen, King). Hearts and diamonds are red;
clubs and spades are black. Face cards = Jack, Queen, King (3 per suit, 12 total).
Work through the tier your teacher assigns.
\end{scenariobox}

\vspace{0.1in}

% ── Tier R ────────────────────────────────────────────────────────────────────
\begin{tcolorbox}[colback=goldbg, colframe=goldacc, boxrule=1pt, arc=2mm,
    left=4mm, right=4mm, top=3mm, bottom=3mm,
    title={\bfseries Tier R \enspace Remediate --- Building Sample Spaces}]
\small

One card is drawn at random from a standard 52-card deck.

\begin{enumerate}[leftmargin=1.5em, itemsep=10pt]
  \item How many total outcomes are in the sample space? \ans{52}

  \item Calculate each probability. Show work as a fraction.
        \vspace{4pt}
        \renewcommand{\arraystretch}{1.8}
        \begin{tabularx}{\linewidth}{@{} l X @{}}
        Event & $P(E) =$ \\
        \hline
        Drawing a heart & \ans{$13/52 = 1/4 = 0.25$} \\
        Drawing a red card & \ans{$26/52 = 1/2 = 0.5$} \\
        Drawing a face card & \ans{$12/52 = 3/13 \approx 0.231$} \\
        Drawing the Ace of Spades & \ans{$1/52 \approx 0.019$} \\
        \end{tabularx}

  \item Apply the complement rule:
        \begin{enumerate}[label=(\alph*), itemsep=6pt]
          \item $P(\text{not a heart}) = $ \ans{$1 - 1/4 = 3/4 = 0.75$}
          \item $P(\text{not a face card}) = $ \ans{$1 - 3/13 = 10/13 \approx 0.769$}
        \end{enumerate}

  \item Is it possible to draw a card with probability 0? Probability 1? Explain.\\
        \ansline{Yes: $P = 0$ for an impossible draw (a ``7 of Stars''); $P = 1$ for a certain one (any card, $52/52$).}
\end{enumerate}

\end{tcolorbox}

\vspace{0.1in}

% ── Tier A ────────────────────────────────────────────────────────────────────
\begin{tcolorbox}[colback=greenbg, colframe=greenacc, boxrule=1pt, arc=2mm,
    left=4mm, right=4mm, top=3mm, bottom=3mm,
    title={\bfseries Tier A \enspace Approaching Proficiency --- Probability and Interpretation}]
\small

One card is drawn at random from a standard 52-card deck.

\begin{enumerate}[leftmargin=1.5em, itemsep=10pt]
  \item Calculate and interpret each probability in the context of the card draw.
        \vspace{4pt}
        \renewcommand{\arraystretch}{1.8}
        \begin{tabularx}{\linewidth}{@{} l p{2.5cm} X @{}}
        Event & $P(E)$ & Long-run interpretation \\
        \hline
        Drawing a club & \ans{$13/52 = 1/4$} & \ansline{Over many draws, about 25\% are clubs.} \\
        Drawing a King & \ans{$4/52 = 1/13$} & \ansline{Over many draws, about 1 in 13 is a King.} \\
        Drawing a red face card & \ans{$6/52 = 3/26$} & \ansline{Over many draws, about 11.5\% are red face cards.} \\
        \end{tabularx}

  \item A friend says: ``I drew a heart, so the probability of drawing a heart is 1.''
        Explain the error in this reasoning.\\
        \ansline{One draw is not a long-run proportion; $P(\text{heart})$ is still $13/52 = 1/4$.}

  \item Use the complement rule to find:
        \begin{enumerate}[label=(\alph*), itemsep=6pt]
          \item $P(\text{not a King}) = $ \ans{$1 - 4/52 = 12/13 \approx 0.923$}
          \item $P(\text{not a red face card}) = $ \ans{$1 - 6/52 = 23/26 \approx 0.885$}
        \end{enumerate}

  \item Use a \textbf{tree diagram} to list the sample space for drawing the suit and
        color of one card. (Hint: there are 2 colors and 4 suits; which suits go with
        which color?) Label each branch and compute $P(\text{red})$ from the tree.
        \par\noindent\begin{minipage}[t][50pt][t]{\linewidth}\ans{Red $\to$ Hearts, Diamonds;\ Black $\to$ Clubs, Spades. Each suit branch has probability 1/4.}\\
        \ans{$P(\text{red}) = P(\text{heart}) + P(\text{diamond}) = 1/4 + 1/4 = 1/2$}\end{minipage}
\end{enumerate}

\end{tcolorbox}

\vspace{0.1in}

% ── Tier E ────────────────────────────────────────────────────────────────────
\begin{tcolorbox}[colback=sky, colframe=navy, boxrule=1pt, arc=2mm,
    left=4mm, right=4mm, top=3mm, bottom=3mm,
    title={\bfseries Tier E \enspace Extension --- Multi-Stage Sample Spaces}]
\small

A spinner has 10 equally likely sections: sections 1--5 are red, sections 6--8 are
blue, and sections 9--10 are green. The spinner is spun \textbf{twice}.

\begin{enumerate}[leftmargin=1.5em, itemsep=10pt]
  \item How many outcomes are in the sample space for two spins? \ans{$10 \times 10 = 100$}
  \item Use a tree diagram or organized list to count the number of outcomes in which
        \textbf{both spins land on red}.
        \par\noindent\begin{minipage}[t][40pt][t]{\linewidth}\ans{Each spin has 5 red sections (1--5): $5 \times 5 = 25$ outcomes.}\end{minipage}
  \item Calculate $P(\text{both red}) = $ \ans{$25/100 = 0.25$}
  \item Calculate $P(\text{at least one blue spin}) $ using the complement rule:\\
        Hint: first find $P(\text{no blue on either spin})$.
        \par\noindent\begin{minipage}[t][40pt][t]{\linewidth}\ans{Not blue: 7 of 10 sections. $P(\text{no blue}) = (7/10)^2 = 0.49$.}\\
        \ans{$P(\text{at least one blue}) = 1 - 0.49 = 0.51$}\end{minipage}
  \item Interpret $P(\text{at least one blue spin})$ in one sentence using the phrase
        ``in the long run.''\\
        \ansline{In the long run, about 51\% of pairs of spins have at least one blue spin.}
  \item (Challenge) If you know that the first spin was red, does the probability that
        the second spin is also red change? Why or why not?\\
        \ansline{No. The spins are independent, so $P(\text{2nd red}) = 5/10 = 0.5$ regardless.}
\end{enumerate}

\end{tcolorbox}

\end{document}
